\subsection{Strong components without a small-category census}

\begin{lemma}[Proved: the categorical minimal ideal]
\label{CP2:A-high-Rees}
Let $S$ be a finite strongly connected category.  There are a
finite group $G$, positive row and column numbers $r_i,c_i$
at each object, and a least nonempty ideal $K$, with
\[
 |K(i,j)|=r_j|G|c_i,\qquad
 |K|=|G|\Bigl(\sum_i r_i\Bigr)\Bigl(\sum_i c_i\Bigr).
\]
If every $r_i=1$, or every $c_i=1$, then $BS\simeq BG$.
If $1_i\in K$, object $i$ is a retract of every object.
When $G=1$, put $A=\mathbb QS/\mathbb QK$.  There are
left and right $A$-modules $V,W$, of dimensions
$\sum(r_i-1)$ and $\sum(c_i-1)$, such that
\[
 H_j(BS;\mathbb Q)=\operatorname{Tor}_{j-2}^{A}(W,V)
 \qquad(j\geq2).
\]
\end{lemma}
\begin{proof}
Fix an object $a$ and an idempotent $e$ in the least ideal
of its endomorphism monoid.  Its corner $G=eS(a,a)e$ is a
group by the finite-monoid argument in
Lemma~\ref{CP2:A-small-minimal-ideal}.  Set
$U_j=S(a,j)e$, $V_i=eS(i,a)$, and $K(i,j)=U_jV_i$.
Every nonempty ideal of $S$ meets $S(a,a)$ by composition
with arrows to and from $a$, and hence contains $e$.
It therefore contains all these $K(i,j)$.  They themselves
form an ideal, proving minimality.

The right $G$ action on $U_j$ and the left action on $V_i$
are free: composing an equation $ug=u$ with any
$v\in V_j$ gives $(vu)g=vu$ inside $G$.  Multiplication
induces a bijection $U_j\times_G V_i\to K(i,j)$.
For injectivity, if $uv=u'v'$, choose $w\in U_i$ and
adjust it by $G$ so that $vw=e$.  Then $u=u'g$ with
$g=v'w\in G$; composing the original equation with a
left inverse of $u'$ gives $gv=v'$.  This is exactly
the balanced equivalence relation.  Choose orbit
representatives.  The resulting coordinates are a row,
a group element, and a column, and composition multiplies
the group coordinates with the intervening fixed element
of $G$.  This gives both cardinality formulas.

Adjoin the retract object represented by $e$ in the
idempotent completion.  It is a retract of $a$, so its
addition preserves the classifying space by
Lemma~\ref{CP2:A-retract-deletion}.  If every row number
is one, the comma category from its group corner to each
object is the free transitive $G$ action on $U_i$.
It has exactly one arrow between any ordered pair of
objects and is contractible.  Quillen's Theorem~A gives
$BG\simeq BS$; the column case is dual.
If $1_i=uv\in K(i,i)$, then $vu$ is an idempotent in
$G$, hence $e$.  For any $j$ choose $u_j,v_j$ with
$v_ju_j=e$.  The maps $u_jv:i\to j$ and $uv_j:j\to i$
compose to $1_i$, proving the retract assertion.

Now suppose $G=1$, and write $R=\mathbb QS$.
The projectives $Re,eR$ augment onto the constant left
and right modules.  Their kernels, now denoted $V,W$,
have the stated dimensions.  The ideal acts constantly
on each row or column set and thus annihilates these
augmentation kernels.  Multiplication is a bimodule
isomorphism $Re\otimes_{\mathbb Q}eR\to\mathbb QK$,
as can be checked on the row-column basis.  In particular
$\mathbb QK$ is projective on both sides.
Dimension-shifting through the two augmentation sequences
gives $H_j(BS;\mathbb Q)=\operatorname{Tor}_{j-2}^R(W,V)$.
At $j=2$ the last map is into $eV=0$, so this includes
ordinary tensor product.  Finally the projective left
resolution $0\to\mathbb QK\to R\to A\to0$ gives
$W\otimes_R\mathbb QK=We\otimes eR=0$ and
$\operatorname{Tor}_{q}^R(W,A)=0$ for $q>0$.
The change-of-rings spectral sequence has only its
zero row and identifies Tor over $R$ with Tor over
$A$.  This proves the formula.
\end{proof}

\begin{lemma}[Proved: two explicit one-sided collapses]
\label{CP2:A-high-one-sided}
The following operations preserve the actual homotopy type.

\begin{enumerate}
\item Suppose $u:b\to a$, $v:a\to b$, and $vu=e\ne1_b$
is idempotent.  If $e$ and every unit of $\operatorname{End}(b)$
fix every incoming arrow from a different object, one may
delete $b$.  There is an outgoing dual.  The unit condition
holds automatically if $e$ is a zero of its corner.
\item Let a strong component $S$ have trivial minimal group
and one minimal row per object.  It may be replaced by one
trivial object if it is a source component.  The new outgoing
stalks are the actual right-act orbit sets.  More generally,
the replacement is allowed if all external incoming arrows
are fixed by local minimal idempotents.  In that case no
hom-set increases.  There are column/sink/exit duals.
\end{enumerate}
\end{lemma}
\begin{proof}
For the first assertion, let $D=C\setminus\{b\}$.
The nonunits of the finite corner at $b$ form an ideal:
if a product is a unit, each factor has a one-sided
inverse and finiteness makes it a unit.  Define
$F:C\to D$ by sending $b$ to $a$, fixing $D$, and
\[
 \begin{aligned}
 F(h)&=1_a\ (h\text{ a unit}),& F(h)&=uhv\ (h\text{ a nonunit}),\\
 F(g)&=ug\ (g:d\to b),& F(f)&=fv\ (f:b\to d).
 \end{aligned}
\]
Identities are preserved.  For nonunits $h,k$, the
equality $ekv=kv$ gives $uhvukv=uhkv$.  For incoming
$g$ it gives $uhvug=uhg$; for outgoing $f$ the equality
$ehv=hv$ gives $fvuhv=fhv$.  Unit factors use $hg=g$
and $hv=v$.  A passage through $b$ between retained
objects gives $fvug=feg=fg$.  In the other order,
$gf:b\to b$ is a nonunit, since $egf=gf$ and a unit
would force $e=1_b$; hence $F(g)F(f)=ugfv=F(gf)$.
These check every pair involving $b$; retained pairs
are unchanged.  Thus $F$ is a functor.  Its composite
with the inclusion is the identity on $D$.  The
components $v$ at $b$ and identities elsewhere give
a natural transformation from the other composite to
$\operatorname{id}_C$, by the same equations.  If $e$
is a zero, $hg=heg=eg=g$ for every unit $h$.

For the second assertion choose $a\in S$ and a minimal
idempotent $e$ there.  There is a unique arrow
$u_i\in S(a,i)e$ at each object, and for every
$f:i\to j$ one has $fu_i=u_j$.  Choose
$v_i\in eS(i,a)$; then $v_iu_i=e$, and put
$e_i=u_iv_i$.  The $u_i$ contract $BS$ by a natural
transformation from the constant functor at $a$.
For a right act $T$, the category of elements has
objects $(i,t)$ and arrows $f:(i,t)\to(j,t')$ with
$t=t'f$.  The map $(i,t)\mapsto tu_i\in T(a)e$
is constant on arrows, and its section sends $z$ to
$(a,z)$.  The $u_i$ give a natural transformation
from section followed by projection to the identity.
Thus the element category is equivalent on classifying
spaces to this discrete orbit set, naturally in $T$.
The two-block pushout proves the source-component
orbit-collapse equivalence.

For the general fixed-entrance assertion, keep the
complement $D$ and introduce one object $s$, with
\[
 \bar C(d,s)=eC(d,a),\qquad \bar C(s,d)=C(a,d)e,
 \qquad \bar C(s,s)=\{1_s\}.
\]
Use the old compositions between retained objects and
through $s$.  No path leaves $S$ and reenters it, since
it is a strong component.  Define $Q$ to send $S$ to
$s$, its arrows to $1_s$, entrances $p:d\to i$ to
$v_ip$, and exits $q:i\to d$ to $qu_i$.
The equations $p=e_ip$ and $fu_i=u_j$ give
$v_jfp=v_ip$, and passage through $S$ gives
$(qu_i)(v_ip)=qp$.  These and associativity verify
every composition.  A section chooses $a$ and the
literal displayed cross arrows.  Its identity at
$s$ is $1_a$.  The identities $v_ae=e$ and $u_a=e$
show that it is indeed a section.  The components
$u_i$ and identities on $D$ give the required natural
transformation to $\operatorname{id}_C$.
Every displayed cross-hom is an image of an old
cross-hom and the internal block becomes one arrow.
This proves the counting assertion as well.
\end{proof}

\begin{lemma}[Proved: nonzero second homology costs three excesses]
\label{CP2:A-high-H2-cost}
If a connected finite category has $k$ objects and at most
$2k+1$ arrows, then $H_2(BC;\mathbb Q)=0$.
\end{lemma}
\begin{proof}
Delete object retracts.  Each deletion removes at least
three arrows and one object and preserves the bound.
Assume none remains.  A proper strong component on
$t\geq2$ objects has at least $t^2+t$ internal arrows:
its minimal ideal has at least $t^2$ and all identities
lie outside it by Lemma~\ref{CP2:A-high-Rees}.
If there are $\ell$ outside objects, contracting the
component in the connected underlying graph leaves at
least $\ell$ edges, and one interface supplies at least
$t$ arrows.  With the outside identities,
\[
 n\geq t^2+t+2\ell+t-1=2k+t^2-1>2k+1.
\]
Thus no such proper component exists.  If the category
is wholly strongly connected, one row or one column
per object gives finite-group homotopy type.  Otherwise
its ideal has at least $(k+1)^2$ arrows, and for
$k\geq2$ its identities give
$n\geq(k+1)^2+k>2k+1$.  The one-object case has order
at most three and rationally acyclic classifying space.

The remaining support is acyclic, with $D+P+R\leq2$
in the notation of the preceding theorem.  If $P=0$,
singleton-support comparison and the flag budget give
$H_2=0$.  If $P>0$, then $D\leq1$ and every corner
algebra is semisimple.  With $R=0$, height one gives
$H_2=0$.  The only other budget is $R=P=1$, $D=0$.
A graph of cycle rank one has either no triangle or
one triangle with bridge attachments.  In the first
case there are no two-simplices.  In the second case
there is only one triangle, and at most one of its
hom-sets is doubled.  Its second boundary is injective:
if an adjacent set is doubled, the boundary component
retaining that set distinguishes the two basis
simplices; otherwise there is just one two-simplex
with nonzero boundary.  Bridges add no two-chains.
This proves the assertion without assuming the
support itself has zero first homology.
\end{proof}

\begin{theorem}[Proved: eight excesses for arbitrary support]
\label{CP2:A-high-all-eight}
Every finite category with $k$ objects and $n\leq2k+7$
arrows has classifying space different from
$\mathbb{CP}^2$.  There is no upper cap on $n$ in this statement.
\end{theorem}
\begin{proof}
Induct on $k$.  For $k=1$ use Corollary~\ref{A:monoids-nine}.
Suppose a connected counterexample has $k>1$ objects.
An object-retract deletion removes at least three
arrows and one object and preserves $n'\leq2k'+7$;
the inductive hypothesis excludes it.  Assume no
object is a retract of another.

A proper strong component on $t\geq3$ objects gives
$n\geq2k+t^2-1\geq2k+8$ by the count just proved.
If the whole category is strongly connected, the
one-row or one-column alternatives have finite-group
type.  Otherwise $n\geq(k+1)^2+k$, forcing $k=2$,
$n=11$, trivial minimal group, nine ideal arrows,
and exactly the two identities outside the ideal.
The quotient algebra in Lemma~\ref{CP2:A-high-Rees}
is $\mathbb Q^2$, so $H_4=0$ by its Tor formula.

It remains to exclude a proper strong pair $S$.
Write $s=|S|$ and $\ell=k-2\geq1$.
It has at least six arrows, and connectedness gives
\[
 s+2\ell+1\leq n\leq2\ell+11,
 \qquad 6\leq s\leq10.
\]
Here a support interface with $S$ has an arrow at each
of its two objects.  The row-column count shows that
$BS$ has finite-group homotopy type: otherwise both
total row and column numbers would be at least three,
requiring nine ideal arrows and two identities.
In particular its second rational cohomology vanishes.

First suppose $S$ is an extreme component; by opposition
make it a source.  If its complement $D$ has $q$ connected
components and nonzero second cohomology, the H2-cost
lemma gives
$|D|\geq2\ell-q+3$.  Every component meets $S$ in
at least two arrows.  Thus
\[
 n\geq s+2\ell+q+3\geq s+2\ell+4.
\]
For $s\geq8$ this contradicts the upper bound.  Both
blocks then have zero second cohomology, and the
two-block lemma excludes the square.  For $s=6,7$,
the minimal group is trivial and each object has one
minimal row and column.  Source orbit collapse in
Lemma~\ref{CP2:A-high-one-sided} saves $s-1$ internal
arrows and at least one interface arrow.  In fact its
orbit set is a surjective image of either nonempty
stalk: if $z\in T(a)e$, then
$(zv_i)u_i=z$.  Its size is at most the smaller of the
two stalk sizes.  The collapse saves at least $s$
arrows and one object, so induction applies.

Now suppose $S$ is nonextreme.  There are outside
objects $a,b$ with $a\to S\to b$, and composition
gives $a\to b$.  The contracted support has at least
$\ell+1$ edges, including two interfaces.  Let
$P_{\rm ext}$ be the sum of excesses $|\operatorname{Hom}|-1$
over all individual nonempty hom-sets between $S$ and
outside objects.  Counting gives
\[
 n\geq s+2\ell+3+P_{\rm ext},\qquad
 P_{\rm ext}\leq8-s,
\]
so $s\leq8$.  For $s=6,7$, the unique minimal
endomorphism $e_i$ at either object is a zero of its
corner.  Consider the four conditions that all incoming
cross arrows, or all outgoing cross arrows, are fixed
by $e_i$.  Each of the at most $s-6$ extra internal
arrows can violate at most two conditions.  An external
violation requires a nonsingleton hom-set and can
violate only the condition at its one incident object.
The number of failed conditions is therefore at most
\[
 2(s-6)+P_{\rm ext}\leq s-4\leq3.
\]
A condition holds.  The first one-sided collapse
deletes that object: its zero ensures the unit
hypothesis, and the minimal arrows to the other
object compose to $e_i$.  At least six arrows
disappear, namely two corner arrows, two internal
cross arrows, and an external entrance and exit.
Induction applies.  For $s=8$ one has
$P_{\rm ext}=0$.  Its minimal group is trivial and
the possible total row/column counts are
$(2,2),(2,3),(3,2)$.  Thus one row or one column occurs
at each object.  The external singleton hom-sets are
fixed, so the second one-sided collapse or its dual
replaces $S$ by a point, saving at least seven arrows
and one object.  Induction again applies.

Every strong-component alternative has been excluded.
The remaining support is acyclic and is excluded by
Theorem~\ref{CP2:A-high-acyclic-eight}.  This completes
the induction.
\end{proof}

\begin{proposition}[Proved: a proper strong pair at the ninth excess]
\label{CP2:A-high-pair-nine}
Suppose $C$ is connected, no object is a retract of
another, and $C$ has a proper strong component on two
objects.  If $n\leq2k+8$, then
$BC\not\simeq\mathbb{CP}^2$.
\end{proposition}
\begin{proof}
Use $S,s,\ell$ as in the preceding proof.  The count
is now $6\leq s\leq11$, with
$s+2\ell+1\leq n\leq2\ell+12$.  For $s\leq10$
the pair has finite-group homotopy type and zero
second rational cohomology.

For an extreme source pair, $s=6,7$ collapses with
at least $s$ arrows saved, reducing to the excess-eight
theorem.  For $8\leq s\leq10$, zero second cohomology
of the complement gives the two-block obstruction.
Otherwise the H2-cost count is
$n\geq s+2\ell+q+3$, forcing
$s=8$, $q=1$, $|D|=2\ell+2$, and exactly two
interface arrows.  All external hom-sets at $S$ are
singleton or empty.  Its one-row or one-column
fixed-interface collapse saves at least seven arrows,
again reducing to the preceding theorem.
For $s=11$, the elementary component count gives
\[
 n\geq11+(2\ell-q)+2q=2\ell+11+q.
\]
Equality forces $q=1$, $|D|=2\ell-1$ and two
interface arrows.  The complement is a trivial-corner
singleton-hom tree.  Both interface arrows end at
one object $w$; it has no outgoing cross arrow in
$D$, since that would create more interface arrows.
The full category on $S\cup\{w\}$ has terminal
object $w$.  Its nerve and $ND$ cover $NC$, meet
at $w$, and are both contractible.  The whole space
is contractible.  Opposites cover extreme sinks.

For a nonextreme pair, the same count becomes
$P_{\rm ext}\leq9-s$ and $s\leq9$.
For $s=6$, and for $s=7$ with $P_{\rm ext}\leq1$,
the four-condition bound is less than four, so
one-sided object deletion saves at least six arrows.
For $s=8,9$ with $P_{\rm ext}=0$, fixed-interface
block collapse saves at least seven.  Each reduces
to Theorem~\ref{CP2:A-high-all-eight}.

The remaining numerical pairs are
$(s,P_{\rm ext})=(7,2),(8,1)$.  Equality holds
in the support count.  All outside corners are
trivial, all outside hom-sets are singleton, the
contracted graph has one triangle $a,S,b$ with
tree attachments, and its only two interfaces with
$S$ are $a\to S$ and $S\to b$.  Delete each outside
tree leaf: it has trivial corner and one singleton
edge to a trivial neighboring corner, so its nerve
is a single interval attached at one endpoint.
This is a homotopy equivalence and removes two
arrows.  We reach a four-object core with sixteen
arrows and a singleton direct arrow $a\to b$.

For $(8,1)$ all parallel excess lies on one interface
side.  The other side has one arrow at each object
of $S$, making $a$ initial or $b$ terminal.  The
core is contractible.  The same holds for $(7,2)$
if both excesses lie on one side.  Otherwise there
is one doubled entrance and one doubled exit.
If any of the four fixed-arrow conditions holds,
one-sided deletion applies as before.  Assume all
four fail.  The unique internal arrow beyond the
four-arrow minimal band and two identities must
then be a cross arrow $t:i\to j$.  It accounts
for failure of the outgoing condition at $i$ and
the incoming condition at $j$.  Therefore the
doubled exterior stalks are $a\to i$ and $j\to b$;
the other two exterior stalks are singletons.

Write $R=\mathbb QS$ and $K$ for its four-arrow
minimal band.  The quotient $A=R/\mathbb QK$ is
the path algebra with two vertices and the one
arrow $t:i\to j$.  Let $V,W$ be the incoming and
outgoing augmentation kernels.  They are killed
by $K$.  At the doubled incoming stalk, its
minimal endomorphism factors through the other
singleton stalk and hence acts constantly; the
other augmentation space is zero.  This proves
$KV=0$, and the dual argument gives $WK=0$.
The projective ideal $\mathbb QK=Re\otimes eR$
and the change-of-rings argument in
Lemma~\ref{CP2:A-high-Rees} identify
$\operatorname{Tor}_2^R(W,V)$ with
$\operatorname{Tor}_2^A(W,V)$.
This last group is zero: a representation
$f:X_i\to X_j$ of the one-arrow quiver has the
explicit projective resolution
\[
 0\to P_j\otimes X_i\to
 (P_i\otimes X_i)\oplus(P_j\otimes X_j)
 \to X\to0,
\]
whose maps at $j$ are $x\mapsto(x,-f(x))$ and
$(x,y)\mapsto f(x)+y$.  The generalized middle-bar
formula in Lemma~\ref{CP2:A-high-linear-Tor}
now gives $H_4$ of the core equal to zero.
Thus even the exceptional cores are excluded
without any sixteen-arrow composition census.
\end{proof}

\begin{corollary}[Proved: every ninth-excess candidate has acyclic support]
\label{CP2:A-high-nine-support}
If $BC\simeq\mathbb{CP}^2$ and $n\leq2k+8$, then
$C$ has acyclic support, no object retract, and
$n=2k+8$.  Consequently a seventeen-arrow candidate
has at most four objects; eighteen- or nineteen-arrow
candidates have at most five; twenty- or twenty-one-arrow
candidates have at most six.  An eighteen-arrow
five-object candidate has acyclic support.
\end{corollary}
\begin{proof}
The equality follows from the excess-eight theorem.
A retract deletion would save at least three arrows
and one object, reducing to that theorem as well.
A proper component on $t\geq3$ objects has
$n\geq2k+t^2-1$.  The only possible case is $t=3$
with equality throughout.  Its ideal has one row
and one column per object and trivial group; it has
twelve internal arrows.  In the contracted support
it is a leaf, with just one interface containing
three singleton arrows.  It is an extreme component
and source-orbit collapse or its dual replaces it
by one point, saving at least thirteen arrows and
two objects.  This reduces to the excess-eight
theorem.

A wholly strong category with $k\geq2$ either has
finite-group type or has $n\geq(k+1)^2+k$.
The latter bound forces $k=2$, trivial minimal
group, nine ideal arrows, and at most three arrows
outside the ideal.  The quotient algebra contains
the two distinct nonzero object idempotents and
is nonlocal of dimension at most three.
Its degree-two Tor cannot be nonzero with bounded
higher Tor, by Lemma~\ref{CP2:A-high-small-Tor}.
The categorical minimal-ideal formula therefore
rules out the projective-plane homology.
A proper strong pair is ruled out by the preceding
proposition.  All nontrivial strong components are
thus excluded, proving acyclicity of the support.
The numerical object bounds are immediate from
$n>2k+7$, and the final assertion uses $18=2(5)+8$.
\end{proof}
